\documentclass[11pt]{article}
\usepackage[T1]{fontenc}
\usepackage[utf8]{inputenc}
\usepackage[english]{babel}
\usepackage{lmodern}
\usepackage{amsmath,amssymb,amsthm}
\usepackage{booktabs,array}
\usepackage[margin=1in]{geometry}
\usepackage{microtype}
\usepackage[colorlinks=true,linkcolor=black,citecolor=blue,urlcolor=blue]{hyperref}
\hypersetup{pdfauthor={Carlo Corti},pdftitle={Exact Enumeration of Adjacent-Sum Permutations: OEIS A180205 through n=28}}
\newtheorem{theorem}{Theorem}[section]
\newtheorem{lemma}[theorem]{Lemma}
\newtheorem{proposition}[theorem]{Proposition}
\newtheorem{corollary}[theorem]{Corollary}
\theoremstyle{definition}
\newtheorem{definition}[theorem]{Definition}
\theoremstyle{remark}
\newtheorem{remark}[theorem]{Remark}
\newcommand{\Znn}{\mathbb Z_{\geq0}}
\newcommand{\An}{\mathcal A_n}
\title{Exact Enumeration of Adjacent-Sum Permutations:\\OEIS A180205 through $n=28$}
\author{Carlo Corti\\Independent researcher\\\texttt{carlo.corti@outlook.com}}
\date{November 2026}
\begin{document}
\maketitle

\begin{abstract}
OEIS A180205 counts permutations of $0,\ldots,n-1$ that are
consecutive sums of two adjacent entries of a nonnegative integer sequence of length
$n+1$. We report an exact enumeration through $n=28$, adding eleven locally
validated values to the historical 17-term baseline, with
$a(28)=1226245772785310491672576$. The unique zero label fixes the entire
realization and separates two outward tails on complementary label sets.
We prove the resulting counting identity and implement its weighted subset
recurrence with a scalar continuation state, parity compression and colexicographic
indexing. An algorithmically distinct inclusion--exclusion enumeration of
repeated-label walks agrees on all 28 values. Both enumerations share the
zero decomposition, scalar transition and ordered tail product on actual
labels. Two coprime moduli and a factorial bound make Chinese remainder
reconstruction exact conditional on correct, complete residues.
A direct interval predicate checks the original definition
exhaustively through $n=8$, and an unpacked arbitrary-precision Python
implementation of the same subset recurrence supplies an additional comparison
through $18$. We give
proofs, intermediate bounds, operation counts and the measured computational
costs, keeping the finite validation domains explicit.
\end{abstract}

\noindent\textbf{2020 Mathematics Subject Classification.}
Primary 05A05; Secondary 05A15, 05A19.

\smallskip
\noindent\textbf{Key words and phrases.}
permutation, adjacent sum, nonnegative realization, subset dynamic programming,
inclusion--exclusion, exact enumeration, OEIS A180205.

\section{Introduction}
\label{sec:intro}

For $n\geq1$, let $\An$ consist of the permutations $p=(p_1,\ldots,p_n)$
of $\{0,\ldots,n-1\}$ for which there exists
$x=(x_0,\ldots,x_n)\in\Znn^{n+1}$ satisfying
\begin{equation}
 p_i=x_{i-1}+x_i\qquad(1\leq i\leq n).
 \label{eq:definition}
\end{equation}
Write $a(n)=|\An|$, with initial index $1$, as in OEIS
A180205~\cite{hardin}. The witness entries may repeat and may be zero.
We count ordered permutations, without identifying reversal, rotation,
complementation or inverse permutations. For example, $(1,0,2)$ is
realized by $(1,0,0,2)$, whereas $(0,2,1)$ has no nonnegative realization.

The historical baseline used here has 17 terms, ending in
$a(17)=445393216256$, and retains R.~H. Hardin's attribution.
The entry gives the author date 15 August 2010, and its OEIS revision history
already includes all 17 terms in revision 1, dated 27 August 2010.
The official OEIS export dated 16 November 2026 agrees with that
baseline~\cite{hardin}. This identifies the comparison range, rather than
asserting that no other results exist. In particular, an earlier preliminary
assessment attributed to Meta AI already listed candidates
$4783835680768$ and $54372476838912$ for indices 18 and 19.
Its original executable calculations and raw logs were not supplied.
The present work does not claim that these numbers first appeared in the
completed runs: their validation here rests on the later recorded
enumerations and comparisons.

The adjacent-window family includes A180206--A180212, with
window lengths three through nine, respectively~\cite{a180206,a180207,a180208,a180209,a180210,a180211,a180212}.
The author's earlier four-window study~\cite{four} uses a zero separator,
weighted subset counting and repeated-label inclusion--exclusion. Its
continuation state is a triple; the present two-window problem has a scalar
state rather than a tuple of the preceding realization entries, permitting
the parity and storage reductions developed below.
That study provides methodological context, not numerical values for A180205.

Inverse adjacent sums are also described in the first-sums composition
entries A391643, A391645 and A391621~\cite{a391643,a391645,a391621}. They use an initial
integer and partial alternating sums to distinguish uniqueness, infeasibility
and witness multiplicity. A391681 counts compositions with parts greater
than one that have no strictly positive realization, while A390567 counts
compositions into positive parts with distinct first sums~\cite{a391681,a390567}.
Here a composition has positive parts, including the empty composition
at total zero; a nonnegative realization is a different condition.
These objects are indexed by composition total or
rank and do not give the distribution by the prescribed label set
$\{0,\ldots,n-1\}$. Their aggregate formulas are therefore not used to
obtain terms of A180205. Another apparent correlation is only numerical:
A131529 counts permutations with distinct adjacent differences~\cite{differences}.
It agrees on the first six values, but has $788$ at index $7$, compared
with $a(7)=768$.

The contribution is a proved scalar-tail enumeration with exact arithmetic,
a complete validated prefix through $28$, and an explicit account of what
each validation method checks, including their shared premises in
Section~\ref{sec:validation}. Scanning all $n!$ permutations is replaced
by weighted aggregation over subsets. Section~\ref{sec:criterion} gives
a direct feasibility test, Sections~\ref{sec:tails}--\ref{sec:producer}
prove and describe the producer, and Section~\ref{sec:ie} gives the separate
inclusion--exclusion count. Sections~\ref{sec:validation}--\ref{sec:results}
state the recorded coverage, results and costs.

\section{A direct feasibility criterion}
\label{sec:criterion}

The following test applies to any nonempty nonnegative integer word $p$,
without locating a zero. Put
\begin{equation}
 c_0=0,\qquad c_j=p_j-c_{j-1}\quad(1\leq j\leq n).
 \label{eq:c}
\end{equation}
For $x_0=t$, successive substitution in~\eqref{eq:definition} gives
\begin{equation}
 x_j=(-1)^j t+c_j.
 \label{eq:intervalwitness}
\end{equation}
Define the integer bounds
\begin{equation}
 L=\max\bigl(\{0\}\cup\{-c_j:1\leq j\leq n,\ j\text{ even}\}\bigr),
 \qquad H=\min_{\substack{1\leq j\leq n\\j\text{ odd}}}c_j.
 \label{eq:interval}
\end{equation}

\begin{proposition}[Interval criterion]
\label{prop:interval}
The word $p$ has a nonnegative integer realization if and only if $L\leq H$.
When feasible, it has exactly $H-L+1$ realizations.
\end{proposition}
\begin{proof}
The recurrence forces~\eqref{eq:intervalwitness} for every choice of $t$.
The constraint at an odd index is $t\leq c_j$, and at an even index it is
$t\geq-c_j$; also $t\geq0$. Thus all entries are nonnegative exactly for
the integers in $[L,H]$. Conversely, every such $t$ gives
$x_{j-1}+x_j=c_{j-1}+c_j=p_j$, proving sufficiency and the multiplicity.
The upper bound is defined even for $n=1$, since $j=1$ is odd.
\end{proof}

This is a linear-time test using integer arithmetic. The alternating-sum
parameterization is established first-sums methodology~\cite{a391643,a391645,a391621};
it is rederived here to specify a validation predicate that does not use
the zero decomposition or the subset recurrence.

\section{Zero decomposition and structural bounds}
\label{sec:tails}

\begin{lemma}[Unique realization]
\label{lem:unique}
Every $p\in\An$ has exactly one nonnegative integer realization.
\end{lemma}
\begin{proof}
Let $p_j=0$. Equation~\eqref{eq:definition} and nonnegativity force
$x_{j-1}=x_j=0$. Moving right, $x_i=p_i-x_{i-1}$ determines the
remaining entries; moving left, $x_{i-1}=p_i-x_i$ does the same.
Every entry is therefore forced, and existence is part of membership in
$\An$. For $n=1$, the witness is $(0,0)$. In particular, feasibility
for a permutation forces $L=H$ in Proposition~\ref{prop:interval}.
\end{proof}

\begin{definition}
An outward tail on a finite set $S$ of positive integer labels is an ordering
$(v_1,\ldots,v_h)$ of $S$ such that the scalar walk
\begin{equation}
 s_0=0,\qquad s_i=v_i-s_{i-1}\geq0\quad(1\leq i\leq h)
 \label{eq:tail}
\end{equation}
remains nonnegative. Let $F(S)$ count these orderings, with
$F(\varnothing)=1$.
\end{definition}

\begin{theorem}[Ordered zero decomposition]
\label{thm:split}
For $n\geq1$, with $U=\{1,\ldots,n-1\}$,
\begin{equation}
 a(n)=\sum_{S\subseteq U}F(S)F(U\setminus S).
 \label{eq:split}
\end{equation}
\end{theorem}
\begin{proof}
Let $p\in\An$, with its unique realization from Lemma~\ref{lem:unique}.
For $p_j=0$, read its left labels outward as
$(p_{j-1},\ldots,p_1)$, starting from $x_{j-1}=0$, and its right labels
as $(p_{j+1},\ldots,p_n)$, starting from $x_j=0$.
Both obey~\eqref{eq:tail}; their label sets are complementary subsets of $U$.

Conversely, take a left tail $\ell=(\ell_1,\ldots,\ell_h)$ on $S$ and a
right tail $r=(r_1,\ldots,r_{n-1-h})$ on $U\setminus S$. If their
generated scalar states are $s_1,\ldots,s_h$ and
$t_1,\ldots,t_{n-1-h}$, form the witness
\[
 (s_h,\ldots,s_1,0,0,t_1,\ldots,t_{n-1-h}).
\]
Its adjacent sums are precisely
$(\ell_h,\ldots,\ell_1,0,r_1,\ldots,r_{n-1-h})$.
The two label sets are disjoint and, together with zero, exhaust
$\{0,\ldots,n-1\}$. This is an admissible permutation, whose zero position recovers both tails
and their sets. The construction is a bijection, including two empty tails
when $n=1$. The split has an ordered left and right side: there is no
division by two and no additional binomial factor.
\end{proof}

Actual labels matter. The valid orders on $\{1,2,3\}$ are
$(1,2,3)$, $(1,3,2)$ and $(2,3,1)$, whereas on $\{1,2,4\}$ the only
valid order is $(1,2,4)$. Thus $F(S)$ cannot be replaced by a function
of $|S|$ or by a rank-standardized set.

\begin{proposition}[Positivity, reversal and growth]
\label{prop:growth}
Every finite positive-label set has $F(S)\geq1$. Also, for $n>1$, $a(n)$
is even, and for every $n\geq1$,
\begin{equation}
 a(n+1)\geq2a(n),\qquad 2^{n-1}\leq a(n)\leq n!.
 \label{eq:bounds}
\end{equation}
\end{proposition}
\begin{proof}
For the increasing order $v_1<\cdots<v_h$, the first state is
$s_1=v_1$. If $0\leq s_{i-1}\leq v_{i-1}<v_i$, then
$s_i=v_i-s_{i-1}$ lies in $[0,v_i]$. Induction proves that this order
is a valid tail. The empty set has its single empty ordering.

Reversing $p$ and $x$ preserves~\eqref{eq:definition}. A permutation of
length greater than one cannot equal its reversal, since its first and last
labels differ. Reversal is therefore a fixed-point-free involution on $\An$.

In any realization of a member of $\An$, each endpoint is at most $n-1$,
as follows from its adjacent sum. Prepending the new label $n$ is realized
by prepending $n-x_0\geq0$ to $x$; appending $n$ is realized by appending
$n-x_n\geq0$. Both maps from $\An$ to $\mathcal A_{n+1}$ are injective
by deletion of the new label. Their images are disjoint: a word of length
$n+1\geq2$ cannot have its unique label $n$ at both ends.
This proves the growth inequality, including $n=1$. Iteration from
$a(1)=1$ gives the lower bound. The upper bound counts all permutations.
\end{proof}

The analogous reversal and end-extension properties occur in the earlier
four-window work~\cite{four}; they are not claimed here as new general
phenomena. Positivity also rules out discarding entire subsets as impossible.

\section{Exact subset dynamic programming}
\label{sec:producer}

\subsection{Weighted scalar states}

Fix a maximum index $N$ and put $M=N-1$. Throughout, $m\leq M$ denotes
the size of a smaller label universe $\{1,\ldots,m\}$. The C program
evaluating the subset recurrence is called the \emph{producer}. For
$S\subseteq\{1,\ldots,M\}$, let $W(S,s)$ count tails on $S$ ending at
state $s$. Initialize $W(\varnothing,0)=1$, with all other empty-set
states zero. Extend $W$ by zero outside $0\leq s\leq M$. Removing a
tail's last label gives the pull recurrence
\begin{equation}
 W(S,s)=\sum_{\substack{v\in S\\v\geq s}}W(S\setminus\{v\},v-s),
 \qquad F(S)=\sum_{s=0}^{M}W(S,s).
 \label{eq:dp}
\end{equation}

\begin{proposition}[Completeness and parity]
\label{prop:dp}
Recurrence~\eqref{eq:dp}, evaluated in increasing $|S|$, counts every
tail exactly once and terminates. Its nonzero states satisfy
\begin{equation}
 0\leq s\leq\max S,\qquad s\equiv\sum_{v\in S}v\pmod2
 \label{eq:parity}
\end{equation}
for nonempty $S$.
\end{proposition}
\begin{proof}
The empty tail is the base case. A valid nonempty tail has a unique last
label $v$, whose preceding state is $v-s\geq0$. Removing it gives the
corresponding term of~\eqref{eq:dp}. Conversely, each counted predecessor
appends $v$ to give a valid tail ending at $s$. Induction on $|S|$ proves
the recurrence, and each dependency decreases $|S|$, proving termination.
The final state is $v$ minus a nonnegative predecessor, so is at most $v$.
Finally, $v-s_{\rm old}\equiv v+s_{\rm old}\pmod2$, proving parity by
induction from zero.
\end{proof}

These are integer multiplicities, not Boolean reachability indicators.
For example, tails $(1,2,3)$ and $(2,3,1)$ end at states 2 and 0,
respectively, while $(1,3,2)$ also ends at 0. The two histories ending at
0 must both be retained.

An ambient label not in $S$ has no role in~\eqref{eq:dp}. Consequently,
a single table of all $F(S)$ for the maximum $M$ supplies every
$a(n)$ with $1\leq n\leq N$ by~\eqref{eq:split} on
$\{1,\ldots,n-1\}$. For $n>1$, the producer pairs $S$ with its complement,
iterates only the subsets omitting the largest label, and doubles the sum.
This is an evaluation of the full ordered sum, rather than a quotient of
the objects counted.

\subsection{Colex layout and parallel ownership}

Represent label $v$ by bit $v-1$. Among equal-cardinality subsets,
colexicographic order compares the largest bit at which they differ:
the set containing that bit is later. Equivalently, compare the increasing
bit lists from their last elements backwards. Let
$b_1<\cdots<b_h$ be the zero-based set bits of a mask. Its zero-based
colexicographic rank among $h$-subsets is
\begin{equation}
 R_h(S)=\sum_{i=1}^{h}\binom{b_i}{i}.
 \label{eq:rank}
\end{equation}
This is the standard binomial ranking formula, also given as equation (1)
of Kruchinin et al.~\cite{colex}. One proof groups subsets by their largest
bit $b_h$: the $\binom{b_h}{h}$ preceding subsets have all bits smaller
than $b_h$, and ranking the remaining $h-1$ bits gives the rest by induction.
The empty subset has rank zero; $\binom bi=0$ for $b<i$.
Deleting $b_j$ therefore gives
\begin{equation}
 R_{h-1}(S\setminus\{b_j+1\})=
 \sum_{i<j}\binom{b_i}{i}+\sum_{i>j}\binom{b_i}{i-1}.
 \label{eq:delete}
\end{equation}
Prefix and suffix binomial sums compute all $h$ deletion ranks in $O(h)$
work per mask. Colex is solely a storage index; it does not change label values.
No approximate unranking or external benchmark constant is used.

For $\pi(S)=(\sum_{v\in S}v)\bmod2$, store the value at state
$s=\pi(S)+2j$ in row cell $j=(s-\pi(S))/2=\lfloor s/2\rfloor$,
using a uniform width $w=\lfloor M/2\rfloor+1$.
Padding beyond $M$ is zero. For a contribution from label $v$, the predecessor
state $v-s$ has parity $\pi(S\setminus\{v\})$, so its cell is
\[
 \frac{v-s-\pi(S\setminus\{v\})}{2}
 =\left\lfloor\frac{v-s}{2}\right\rfloor.
\]
The predecessor pulls stay within states $0,\ldots,M$; row summation may
also scan zero padding.
Two cardinality layers suffice in memory. A separate array stores all
$2^M$ values of $F$, at offset
$\sum_{k<h}\binom Mk+R_h(S)$.

For each cardinality, workers own disjoint output rows, read the previous
layer and binomial table, and write their own $F(S)$. A barrier separates
successive layers. Thus no atomic addition to a shared count cell is needed.
The scalar mathematical recurrence and these indexing formulas specify a
reimplementation independently of the checkpoint format.

\subsection{Arithmetic and work bounds}
\label{sec:arithmetic}

For $h=|S|$, each state count and $F(S)$ are at most $h!$.
For complementary subsets of $m$ labels,
$F(S)F(U\setminus S)\leq h!(m-h)!\leq m!$.
Every partial nonnegative sum in the recurrence, row total or ordered
convolution is bounded by the corresponding complete count; in particular
the convolution is at most $n!$. The doubled half-sum also has this bound,
since it consists of disjoint complementary pairs.

The C17 producer uses the GNU unsigned 128-bit integer
extension~\cite{gcc}, with checked additions and multiplications.
Its implemented horizon is $N\leq30$, and
\[
 30!=265252859812191058636308480000000<2^{128}.
\]
The factorial bounds cover intermediates, not merely the anticipated final
answer. Masks use 32 bits; ranks, offsets and work counters use 64 bits,
with checked allocation and file-offset domains. The capacity through 30
is a software limit, not a computed prefix through 30.

For 16-byte cells, the principal resident-array requirement is
\begin{equation}
 B(M)=16\,2^M+16w\max_{1\leq h\leq M}
 \left(\binom M{h-1}+\binom Mh\right),
 \label{eq:memory}
\end{equation}
for $M=N-1\geq1$, with $M=0$ handled separately. This includes $F$ and the two
simultaneously resident layers, not process RSS or the disk layout.

For $M\geq3$, the exact number of predecessor-cell reads in the dense
pull loops is
\begin{equation}
 D(M)=2^{M-3}M(M+3).
 \label{eq:D}
\end{equation}
Indeed, fix an included label $v$. Among its $2^{M-1}$ masks, parity is
balanced: since $M\geq3$, at least one of the odd labels 1 and 3 differs
from $v$, and toggling its membership toggles parity. The mean number of
$s\in[0,v]$ with the prescribed parity is
$(v+1)/2$. Summing over $v=1,\ldots,M$ gives~\eqref{eq:D}.
The small cases are $D(0)=0$, $D(1)=1$, $D(2)=5$.
Rank construction, row clearing and summation, complementary products,
and persistence are additional work. The arithmetic work is
$O(M^2 2^M)$; the principal resident-cell count is
$O(\sqrt{M}\,2^M)$ by the central binomial bound.

\section{Repeated-label inclusion--exclusion and exact CRT}
\label{sec:ie}

\subsection{A different aggregation}

For $T\subseteq\{1,\ldots,m\}$, allow labels to repeat. Let
$K_h(T,s)$ count nonnegative walks of length $h$ starting at zero,
using labels in $T$, and ending at $s$. Initialize $K_0(T,0)=1$ and
all other states zero. The recurrence is
\begin{equation}
 K_{h+1}(T,t)=\sum_{\substack{v\in T\\v\geq t}}K_h(T,v-t),
 \qquad L_h(T)=\sum_{s\geq0}K_h(T,s).
 \label{eq:walk}
\end{equation}
Every state lies in $[0,\max T]$: after a transition using $v$, it is
$v$ minus a nonnegative predecessor. For $T=\varnothing$, only the empty
walk survives. Thus the displayed state sum is finite and independent of
the ambient universe; whenever $\max T\leq m\leq M$, summing through
$m$ or through $M$ gives the same value. Set $\max\varnothing=0$.
Put $Q_m(T)=\sum_{k=0}^{m}L_k(T)L_{m-k}(T)$; the subscripted $Q_m$
is a walk-pair count, distinct from the modulus $Q$ below.

\begin{theorem}[Inclusion--exclusion identity]
\label{thm:ie}
For every $m\geq0$,
\begin{equation}
 a(m+1)=\sum_{T\subseteq\{1,\ldots,m\}}
 (-1)^{m-|T|}Q_m(T).
 \label{eq:ie}
\end{equation}
\end{theorem}
\begin{proof}
The product for a given $k$ counts an ordered pair of repeated-label walks with a total
of $m$ label positions, using only $T$. Fix such a pair, and let $B$ be
the set of labels actually used. Its total weight in~\eqref{eq:ie} is
\[
 \sum_{B\subseteq T\subseteq\{1,\ldots,m\}}(-1)^{m-|T|}
 =\begin{cases}1,&B=\{1,\ldots,m\},\\0,&\text{otherwise}.\end{cases}
\]
Using all $m$ labels in exactly $m$ positions forces every label to appear
once. All sums are finite, so summing by pairs before summing by alphabets
is valid. The surviving pair is therefore a pair of distinct-label tails on
complementary sets. Theorem~\ref{thm:split} proves the identity.
For $m=0$, only two empty tails occur and the value is 1.
\end{proof}

The C++17 validator enumerates allowed alphabets $T$, with small scalar
arrays indexed by walk length and endpoint; it has no used-label mask
inside a walk. It reads no producer $F$ table, colex rank, packed state row
or expected numerical answer. For a maximum $M$, it computes $L_h(T)$
through length $M$ once per alphabet. The same alphabet contributes to
each smaller universe $\{1,\ldots,m\}$ with $m\geq\max T$, using
the sign $(-1)^{m-|T|}$ and $Q_m(T)$.
An alphabet containing a larger label contributes to no such smaller universe.
This gives the full prefix in a single complete $2^M$-alphabet pass.

The implementation loops over every $s=0,\ldots,v$ in the transition
from label $v$, for each of $M$ lengths and two residue channels.
For $M\geq1$, its exact number of these modular additions is
\begin{equation}
 I(M)=2^{M-1}M^2(M+3).
 \label{eq:I}
\end{equation}
Each label belongs to $2^{M-1}$ alphabets and contributes $2M(v+1)$
additions, proving the formula; $I(0)=0$.
This excludes the endpoint totals, $Q_m$ products and sums, and state saves.
It is a different work unit from $D(M)$ and is not a bound for one
checkpoint quantum. For $M\geq3$, $I(M)=4M D(M)$.

\subsection{Residues and reconstruction}

Repeated-label walk counts and signed inclusion--exclusion partial sums
need not obey a factorial bound. Instead, the validator evaluates all
operations modulo
\begin{equation}
 P=2^{61}-1=2305843009213693951,\qquad
 Q=2^{61}-3=2305843009213693949.
 \label{eq:moduli}
\end{equation}
Both are odd and their difference is two, so $\gcd(P,Q)=1$.
No primality assertion about either modulus is required. Residue addition
and subtraction fit in unsigned 64-bit arithmetic, and multiplication of
two reduced residues is smaller than $2^{122}$, fitting in unsigned 128 bits.
Each operation is reduced modulo its channel's modulus.

If the two complete residues for $a(n)$ are $r_P,r_Q$, the separate
arbitrary-precision Python reconstruction computes
\begin{equation}
 z=r_P+P\bigl((r_Q-r_P)P^{-1}\bmod Q\bigr),\qquad 0\leq z<PQ.
 \label{eq:crt}
\end{equation}
Here
\[
 PQ=5316911983139663482391856204266602499>30!.
\]
The bound $0\leq a(n)\leq n!$ therefore makes $z=a(n)$ for $n\leq30$,
\emph{conditional on the residues being correct and complete}.
A reconstructed value above $n!$ is rejected. A value below that bound
alone cannot establish correct residue generation. Complete alphabet
coverage, residue/state agreement and source/build identities are separately
checked before comparison with the producer.

\section{Implementation, evidence and validation boundaries}
\label{sec:validation}

\subsection{Recorded domains and dependence}

Table~\ref{tab:validation} distinguishes complete counts from checks on a
smaller domain. These are retained local executions and comparisons,
not new enumerations performed to prepare the manuscript.

\begin{table}[htbp]
\centering\small
\begin{tabular}{@{}>{\raggedright\arraybackslash}p{0.30\textwidth}
 >{\raggedright\arraybackslash}p{0.15\textwidth}
 >{\raggedright\arraybackslash}p{0.47\textwidth}@{}}
\toprule
Method & Domain checked & What is checked \\
\midrule
C weighted subset DP & $1,\ldots,28$ & Full producer counts, with all smaller prefixes reused \\
C++ repeated-label IE plus Python CRT & $1,\ldots,28$ & Complete recount with different aggregation and modular arithmetic \\
Unpacked scalar Python & $1,\ldots,18$ & Arbitrary-precision subset recurrence without colex or parity packing \\
All permutations with interval predicate & $1,\ldots,8$ & Definition-level membership and complete small totals, without zero splitting \\
Direct orders on each positive-label subset & All 128 subsets of $\{1,\ldots,7\}$ & Every small $F(S)$, including merged-state multiplicities \\
\bottomrule
\end{tabular}
\caption{Actual local validation coverage. The Python scalar method was
qualified through 17 and then applied to the completed run through 18;
its evidence is inherited, not extended, in later runs.}
\label{tab:validation}
\end{table}

The producer and IE share Theorem~\ref{thm:split}, the scalar transition,
the actual label semantics and the product of the two tails. They are
algorithmically distinct counts, not independent proofs of that theorem.
The unpacked Python comparison also shares the zero-tail mathematics,
subset recurrence and parity restriction, while storing full scalar rows.
Its value is a different representation and arithmetic
implementation at its stated indices.

The interval predicate in Proposition~\ref{prop:interval} supplies a
different mathematical test directly from the definition. The local suite
exhaustively tests every permutation through $n=8$, obtaining
$1,2,4,12,44,176,768,4000$. Its initial finite upper bound $H=n$ is
harmless: the first constraint immediately tightens it to
$H\leq p_1\leq n-1$. This check uses neither $F$, parity packing, colex
nor the position of zero. It corroborates the shared reduction on that
small domain, without constituting a direct census through 28.
For the coefficient check, every ordering of each of the 128 subsets of
$\{1,\ldots,7\}$ was tested by~\eqref{eq:tail}, its accepted orders counted,
and that integer compared with the corresponding stored $F(S)$.
The retained validation source and execution summary record this exhaustive
small-subset comparison.

All local numerical checks belong to the OpenAI/Codex-assisted maintenance
workflow under the author's authorization; the C and C++ programs also
share the GCC compiler family. Here algorithmic independence specifies
the separation just described, not a verified third-party replication.
External code-review reports describe further definition-based DP checks
and fault sweeps; manuscript reviewers also report larger definition-based
checks. Complete source and raw execution evidence for those reports were
not all supplied. Those statements remain reviewer-reported evidence and
do not enlarge the local domains in Table~\ref{tab:validation}.

\subsection{Persistence and software qualification}

The producer writes fixed-width little-endian counts and a chained journal
of committed row ranges, binding source, build recipe, executable, maximum
and encoding. Two layers are resident; three rotating layer files preserve
the durable recovery boundary. Data are synchronized before the journal
and control record commit. Recovery checks committed ranges and semantic
state constraints before continuing. Earlier overwritten $W$ layers are
not retained as complete historical layers: their transitions are not
recomputed by this integrity audit. The serialized checkpoint layout is
an implementation detail; mathematical reproduction uses the recurrences
and indexing above. An incomplete run contributes no
validated sequence value. Read-only integrity checking does not recompute
the recurrence, and a digest authenticates retained bytes rather than
proving the mathematical count.

Interrupted-run, explicit-resume, no-clobber, invalid-input, thread-count
and corruption regressions are recorded on bounded known domains.
Current-source AddressSanitizer and UndefinedBehaviorSanitizer checks cover
the prefix through 17 and its audit with leak detection disabled.
Valgrind checks used a separate \texttt{-march=x86-64-v3} build compatible
with the installed Valgrind, and cover 8 and its audit with zero reported
errors or definite leaks, while a possible OpenMP runtime allocation is
disclosed. ThreadSanitizer/Archer qualification is inherited from earlier
instrumented builds with an unchanged parallel kernel; it is not a fresh
clean ThreadSanitizer execution of the final production binary.
These checks support their tested boundaries, not a universal absence of
software defects or instrumented execution through 28.

Original validator sources, their required inputs and execution records
are retained locally. A reproducibility check can reconstruct the final
28-row result from the small stored IE residue/control files and compare
it with the native table without repeating either census. The original
large native intermediate tables have been intentionally discarded after
retaining the counts, completion records, source identities and validation
inputs. A full native-table audit can therefore no longer be replayed on
those original bytes; generating new tables would require a new execution.
The retained journal and residue checks do not recompute those tables.
The retained known-domain regression suite was originally run against its
17-row baseline and must not be interpreted as a fresh test against the
subsequently extended canonical file.

The intended computational package is identified in~\cite{corti}.
Its locators are prospective. Future public-package reproduction is
separate from the retained local execution evidence used here.

\section{Results and measured costs}
\label{sec:results}

\begin{proposition}[Recorded computational result]
\label{prop:result}
The author's runs completed on 16 November 2026 give the values of $a(n)$
in Table~\ref{tab:terms} for $1\leq n\leq28$.
The subset dynamic program and inclusion--exclusion/CRT agree row for row.
In particular,
\[
 a(28)=1226245772785310491672576.
\]
\end{proposition}
\begin{proof}[Computational proof]
Theorem~\ref{thm:split} and Proposition~\ref{prop:dp} establish the
exhaustive ordered producer count; Theorem~\ref{thm:ie} establishes the
second count. The intermediate bounds in Section~\ref{sec:arithmetic}
and the modular/CRT bounds in Section~\ref{sec:ie} cover the implemented
arithmetic. The final completed run, on 16 November 2026, contains all
28 native rows. Its IE control reaches all $2^{27}=134217728$ alphabets.
The producer's journal descriptors partition the ranks
$0,\ldots,\binom{27}{h}-1$ consecutively at each cardinality $h$;
the terminal frontier is beyond the last layer. Colex generation is
bijective on those ranks, so the stated loops cover every subset row.
For $n\geq2$, the complementary-product loop evaluates successive chunks
$[c,\min(c+q,2^{n-2}))$, where $c$ is its current cursor and $q$ the
configured row quantum. Each chunk advances $c$ to its upper endpoint;
the index completes only when that endpoint reaches $2^{n-2}$.
The case $n=1$ is initialized directly to 1.
The journal audit checks increasing cursors for partial chunks and records
index completion, but does not separately check the last chunk's length.
Coverage of that chunk rests on the loop semantics, rather than an
independent journal partition check.
The IE cursor advances consecutive disjoint alphabet ranges from zero
to $2^{27}$, evaluating every alphabet once. The recorded producer reads
are $D(27)=13589544960$, and the recorded IE transition additions are
$I(27)=1467670855680$, agreeing with the exact work identities.
The producer rejects a read counter above $D(27)$; IE rederives its
expected addition count from the cursor. Counter equality corroborates
this traversal; equality alone would not prove coverage, correct arithmetic
or faithful execution. Completion and validation records bind the counts
to the frozen programs. Exact CRT reconstruction agrees
row for row with the native table, and all previous 26 rows agree with
the earlier validated prefix. The supplementary checks and shared premises
have precisely the scopes stated in Section~\ref{sec:validation}.
This conclusion uses the proved recurrences, their implementations and
the recorded executions; it is not an independent arithmetic certificate
derived from the counters.
\end{proof}

\begin{table}[htbp]
\centering\small
\begin{tabular}{@{}rr@{\qquad}rr@{}}
\toprule
$n$ & $a(n)$ & $n$ & $a(n)$ \\
\midrule
1 & 1 & 15 & 4455645952 \\
2 & 2 & 16 & 43468799232 \\
3 & 4 & 17 & 445393216256 \\
4 & 12 & 18 & 4783835680768 \\
5 & 44 & 19 & 54372476838912 \\
6 & 176 & 20 & 653068724224000 \\
7 & 768 & 21 & 8169310794943488 \\
8 & 4000 & 22 & 106298966614843392 \\
9 & 22896 & 23 & 1449278709531697152 \\
10 & 142336 & 24 & 20677932905765044224 \\
11 & 967040 & 25 & 305457397455433007104 \\
12 & 7258368 & 26 & 4667597748152564613120 \\
13 & 57919872 & 27 & 74207848006226125996032 \\
14 & 490253312 & 28 & 1226245772785310491672576 \\
\bottomrule
\end{tabular}
\caption{The validated contiguous prefix. Indices $1$--$17$ reproduce the
historical baseline attributed to Hardin~\cite{hardin}; indices $18$--$28$
are the eleven additional locally validated values. The earlier preliminary
appearance of candidates 18 and 19 is distinguished in Section~\ref{sec:intro}.}
\label{tab:terms}
\end{table}

The canonical b-file contains exactly the 28 indexed data rows, with one
terminal empty line. The native table has the same indexed values.
Every value satisfies Proposition~\ref{prop:growth}. The first index not
computed locally is $29$: neither the bounds nor the available related
sequence data evaluate it from the retained results. No count at $29$ or
$30$ is asserted, and this local endpoint makes no claim about the
absence of external results.

\subsection{Environment and elapsed service}

The recorded host was an AMD Ryzen 9 7940HS with 8 cores and 16 hardware
threads, running Linux Mint 22.3 with x86-64 kernel 6.17.0-42-generic and
64 GB DDR5 installed RAM. The operating system reported
66136649728 bytes of total memory at qualification.
Both final numerical programs requested eight OpenMP workers, with dynamic
teams disabled, placement by cores and binding enabled.
The recorded toolchain is GCC/G++ 13.3.0, Python 3.12.3 and OpenSSL 3.0.13.
The producer's numerical optimization flags were
\texttt{-std=c17 -O3 -march=native -mtune=native -flto}
\texttt{-funroll-loops -fomit-frame-pointer -fopenmp}; the IE build used
\texttt{-std=c++17 -O3 -march=native -fopenmp}. Warning flags were
\texttt{-Wall}, \texttt{-Wextra}, \texttt{-Wpedantic}, \texttt{-Wshadow},
\texttt{-Wconversion}, \texttt{-Wdouble-promotion} and \texttt{-Wformat=2}.
Both programs link against
\texttt{libcrypto}; OpenMP uses
\texttt{libgomp}. These hardware-specific binaries are retained historical
artifacts, not interchangeable with a fresh build.

The recorded numerical options can be expressed by the following templates,
where \texttt{NEW} is an existing parent directory and its
\texttt{prod} and \texttt{ie} children must not already exist:
\begin{small}
\begin{verbatim}
src/a180205 run --root NEW/prod --max 28 \
    --threads 8 --seconds 3000 --quantum 262144
validation/ie run --root NEW/ie --max 28 \
    --threads 8 --seconds 3000 --quantum 256
python3 -B validation/crt.py NEW/ie/cnt.tsv --ie validation/ie
\end{verbatim}
\end{small}
The producer additionally requires the preserved historical 17-term prefix
in its canonical b-file location and verifies its fixed digest after
canonicalizing that prefix. CRT checks that the IE source, build recipe,
executable, residue text and completed control record have the bound
identities. Rebuilt binaries do not replace the original execution identity.

For the final maximum 28, measured producer elapsed service was about
51 seconds and IE service about 1133 seconds. The sum of the four
command-level elapsed wall times, including validation and consolidation,
was about 1185 seconds (19 minutes 45 seconds), excluding operator pauses
between commands. Both numerical engines completed
in their first invocation, without resume. The IE measurement includes its
state-save overhead; $I(M)$ alone is not a wall-time prediction.

Equation~\eqref{eq:memory} gives 11133602048 bytes, or about 10.37 GiB,
for the principal arrays at $M=27$. Adding the implementation's 256 MiB
bookkeeping allowance gives about 10.62 GiB. These are allocation models:
the final run's peak RSS was not measured. A separate earlier synthetic
capacity-30 fixture allocated 45818139392 bytes of principal arrays and
recorded 45821034496 bytes of peak RSS, about 42.67 GiB, using
\texttt{getrusage}. This is a synthetic qualification measurement, not
the run28 footprint or a computation through 30. Units GiB and MiB mean
powers of two.

\section{Conclusion}

The unique zero label turns the inverse adjacent-sum condition into two
ordered scalar walks. Retaining their integer multiplicities on actual
label sets permits a complete subset enumeration, while parity and colex
indexing reduce its memory requirements without changing the objects counted.
Repeated-label inclusion--exclusion provides a second complete aggregation;
coprime-modulus reconstruction is exact within the proved factorial bound.
The resulting contiguous prefix reaches $28$, eleven values beyond the
historical baseline used in this study. Definition-level testing is exhaustive
only on its stated small domain, and both large enumerations retain a shared
mathematical premise. The 28 finite values do not establish an asymptotic law
or an exact recurrence for subsequent terms.

\section*{Data availability}

A local publication package accompanies this monolithic LaTeX manuscript;
its public release is pending. It contains the producer and build recipe,
complete count table and b-file, IE/CRT sources and original residue records,
supplementary Python checks, selected original completion and qualification
records, environment information and a bounded package verifier.
Its indexed compressed evidence includes the original build identities and
inputs required for the documented retained-result checks, without a private
workspace dependency. Original large intermediate tables and obsolete
diagnostic payloads are not retained or required by those checks.
The resulting native-table replay limitation is stated in
Section~\ref{sec:validation}. Enquiries may be addressed to the author
at the email on the title page.
The project reference~\cite{corti} specifies the
intended GitHub repository and Zenodo archive; these are prospective
locators, and the DOI is an author-designated placeholder. No existing public
A180205 deposit is asserted.

\section*{AI use disclosure}

\begingroup\small\setlength{\baselineskip}{11pt}
\noindent\textbf{Mathematical analysis.}
OpenAI Codex generated and revised proof drafts and supported algorithm
selection, arithmetic bounds and source reconciliation during
15--17 November 2026. Exact historical model revisions were not recorded
consistently. An earlier preliminary report
attributed to Meta AI supplied the unadmitted candidates at 18 and 19.
Local methodological notes, including the earlier four-window study, informed
the strategy; an attempted live NotebookLM consultation did not produce a
usable response and is not cited as evidence. Mathematical assertions rest
on the proofs, primary sources and retained artifacts, not on conversational
assertions or model agreement.

\smallskip
\noindent\textbf{Algorithm and code.}
OpenAI Codex generated and revised the C17 producer, C++17 IE validator
and Python CRT driver, and supported validation design and local artifact
checks. The maintainer name Luna in the source headers denotes this
assistant. The author selected the computational objectives and stopping
point and executed the official numerical runs and their validation steps.
These acts are distinct from the assistant-assisted programmatic checks;
no unaided human re-enumeration or individually logged unaided rederivation
of every proof is claimed. Code-review reports
used the tool or vendor designations ChatGPT, Claude, DeepSeek, Gemini,
GLM, Grok, Kimi, Meta AI, MiniMax, Mistral and Qwen; follow-up reports
came from Claude, MiniMax, Meta AI and GLM. Exact backend identities and
versions were not consistently authenticated. A supplied MiniMax review-side
Python implementation was retained, inspected, qualified locally and used
for the scalar comparison through 18. Reported external computations are
distinguished from retained local execution evidence.

\smallskip
\noindent\textbf{Manuscript.}
OpenAI/Codex assisted with drafting, literature checks, internal review,
compilation and artifact verification. Numerical claims were reconciled with
the frozen count tables, b-file, validator sources and original receipts.
The first external manuscript round supplied reports designated ChatGPT,
Claude, DeepSeek, Grok, Meta AI, MiniMax and Qwen. The focused second round
supplied Claude and MiniMax reports. A separate first-round MiniMax Python
check was also executed locally (direct predicate through 8, scalar recurrence
through 10, stored-residue CRT through 28), without enlarging the validation
domains. The reports' findings were
checked against sources and artifacts; claimed executions without full
reproducible evidence remain reported observations. These designations do
not authenticate exact backends. Internal checks and same-family agreement
do not constitute independent third-party replication.

\smallskip
Carlo Corti is the sole author and takes full scientific responsibility for
the mathematical statements, implementation, numerical results, validation
boundaries and conclusions.\par
\endgroup

\clearpage
\begin{thebibliography}{99}

\bibitem{hardin}
R. H. Hardin,
\emph{Number of permutations of 0..(n-1) representable as consecutive
sums of 2 adjacent elements of a sequence of n+1 nonnegative integers},
OEIS A180205, author date 15 August 2010,
\url{https://oeis.org/A180205}. OEIS revision 9, with the historical
17-term data already in revision 1; official OEIS export dated
16 November 2026, consulted November 2026;
\url{https://github.com/oeis/oeisdata}.

\bibitem{a180206}
R. H. Hardin,
\emph{Number of permutations of 0..(n-1) representable as consecutive sums of 3 adjacent elements of a sequence of n+2 nonnegative integers},
OEIS A180206, author date 15 August 2010,
\url{https://oeis.org/A180206}.
Official OEIS export consulted November 2026.

\bibitem{a180207}
R. H. Hardin,
\emph{Number of permutations of 0..(n-1) representable as consecutive sums of 4 adjacent elements of a sequence of n+3 nonnegative integers},
OEIS A180207, author date 15 August 2010,
\url{https://oeis.org/A180207}.
Extension through $n=22$ by Carlo Corti, 13 November 2026.
Official OEIS export consulted November 2026.

\bibitem{a180208}
R. H. Hardin,
\emph{Number of permutations of 0..(n-1) representable as consecutive sums of 5 adjacent elements of a sequence of n+4 nonnegative integers},
OEIS A180208, author date 15 August 2010,
\url{https://oeis.org/A180208}.
Official OEIS export consulted November 2026.

\bibitem{a180209}
R. H. Hardin,
\emph{Number of permutations of 0..(n-1) representable as consecutive sums of 6 adjacent elements of a sequence of n+5 nonnegative integers},
OEIS A180209, author date 15 August 2010,
\url{https://oeis.org/A180209}.
Official OEIS export consulted November 2026.

\bibitem{a180210}
R. H. Hardin,
\emph{Number of permutations of 0..(n-1) representable as consecutive sums of 7 adjacent elements of a sequence of n+6 nonnegative integers},
OEIS A180210, author date 15 August 2010,
\url{https://oeis.org/A180210}.
Official OEIS export consulted November 2026.

\bibitem{a180211}
R. H. Hardin,
\emph{Number of permutations of 0..(n-1) representable as consecutive sums of 8 adjacent elements of a sequence of n+7 nonnegative integers},
OEIS A180211, author date 15 August 2010,
\url{https://oeis.org/A180211}.
Official OEIS export consulted November 2026.

\bibitem{a180212}
R. H. Hardin,
\emph{Number of permutations of 0..(n-1) representable as consecutive sums of 9 adjacent elements of a sequence of n+8 nonnegative integers},
OEIS A180212, author date 15 August 2010,
\url{https://oeis.org/A180212}.
Official OEIS export consulted November 2026.

\bibitem{four}
C. Corti,
\emph{Exact Enumeration of Four-Window Permutations: OEIS A180207 through
$n=22$}, 2026,
\url{https://doi.org/10.5281/zenodo.24194032}.
The DOI identifies the computational package release series; the cited
manuscript is its accompanying article.

\bibitem{a391643}
G. Wiseman,
\emph{Number of integer compositions n that are the first sums of a unique nonnegative sequence},
OEIS A391643, author date 4 January 2026, first index 1,
\url{https://oeis.org/A391643}.
Formula and table contributions by C. Sievers. Official entry consulted November 2026.

\bibitem{a391645}
G. Wiseman,
\emph{Number of integer compositions of n that are not the first sums of any finite nonnegative sequence},
OEIS A391645, author date 5 January 2026, first index 1,
\url{https://oeis.org/A391645}.
Formula and table contributions by C. Sievers. Official entry consulted November 2026.

\bibitem{a391621}
G. Wiseman,
\emph{Number of nonnegative integer sequences whose first sums are the k-th composition in standard order},
OEIS A391621, author date 19 December 2025, first index 1,
\url{https://oeis.org/A391621}.
Table by A. Eldar. Official entry consulted November 2026.

\bibitem{a391681}
G. Wiseman,
\emph{Number of integer compositions of n with all parts > 1 that are not the first sums of any composition},
OEIS A391681, author date 6 January 2026, first index 1,
\url{https://oeis.org/A391681}.
Formula and table contributions by C. Sievers. Official entry consulted November 2026.

\bibitem{a390567}
G. Wiseman,
\emph{Number of integer compositions of n with all distinct first sums},
OEIS A390567, author date 15 November 2025, first index 0,
\url{https://oeis.org/A390567}.
Extension through total 40 by C. Sievers. Official entry consulted November 2026.

\bibitem{differences}
V. Jovovic,
\emph{Number of permutations of $\{1,2,\ldots,n\}$ for which differences
of adjacent numbers are all distinct},
OEIS A131529, 26 August 2007,
\url{https://oeis.org/A131529}.
Extensions by R. J. Mathar and R. H. Hardin; consulted November 2026.

\bibitem{colex}
V. Kruchinin, Y. Shablya, D. Kruchinin and V. Rulevskiy,
Unranking Small Combinations of a Large Set in Co-Lexicographic Order,
\emph{Algorithms} \textbf{15}(2) (2022), article 36,
\url{https://doi.org/10.3390/a15020036}.

\bibitem{gcc}
GNU Project,
\emph{Using the GNU Compiler Collection}, GCC 13.3.0,
Section 6.9, 128-bit Integers,
\url{https://gcc.gnu.org/onlinedocs/gcc-13.3.0/gcc/_005f_005fint128.html},
consulted November 2026.

\bibitem{corti}
C. Corti,
\emph{Exact Enumeration of Adjacent-Sum Permutations: OEIS A180205 through
$n=28$}, GitHub repository and Zenodo archive, 2026.
\url{https://github.com/carcorti/A180205};
\url{https://doi.org/10.5281/zenodo.xxxxxxxx}.

\end{thebibliography}
\end{document}
